% Relativistic Chapter X, §90--92: Rayleigh--Taylor Perturbation Equations (Inviscid Case)
% Relativistic generalization of Chandrasekhar, Chapter X §90--92
% Conventions: see RELATIVISTIC_CONVENTIONS.md

\chapter{The Stability of Superposed Relativistic Fluids: The Rayleigh--Taylor Instability}
\label{ch:rel-10}

%% ================================================================
\section{Introduction}\label{sec:rel-90}

In the non-relativistic theory (Chapter~X, \S\ref{sec:90}), the
Rayleigh--Taylor instability arises when a heavy fluid is supported above a
lighter one in a gravitational field, or---more generally---when a density
interface is subjected to an acceleration directed from the heavy to the
light fluid.  The instability plays a fundamental role in inertial
confinement fusion, atmospheric dynamics, and astrophysical mixing.

In the relativistic regime the same physical mechanism operates, but with
important modifications:
\begin{enumerate}
\item \textbf{Inertia of internal energy.}  The role of the mass density
  $\rho$ is taken over by the relativistic enthalpy per unit volume,
  \begin{equation}
  \label{eq:rel-10-w-def}
  w \;=\; \frac{\varepsilon + p}{c^{2}},
  \end{equation}
  where $\varepsilon$ is the total energy density and $p$ the pressure.
  Because pressure contributes to inertia, the effective ``weight'' of a
  fluid element differs from its rest-mass density.

\item \textbf{Causality constraints.}  Although the Rayleigh--Taylor mode
  is non-oscillatory (the growth rate~$\sigma$ is real), the perturbation
  equations must still respect $v \le c$.  In particular the sound speed
  $c_{\mathrm{s}}^{2} = (\partial p/\partial\varepsilon)_{s}$ must
  satisfy $c_{\mathrm{s}} \le c$, and any dissipative extension must
  employ a causal first-order formalism---such as BDNK
  (Bemfica--Disconzi--Noronha--Kovtun) theory
  \citep{Bemfica2018,Kovtun2019}---rather than parabolic diffusion.

\item \textbf{Astrophysical contexts.}  Relativistic Rayleigh--Taylor
  instability is encountered in:
  \begin{itemize}
  \item \emph{Relativistic jets} from active galactic nuclei and
    micro-quasars, where the jet--cocoon interface is RT-unstable under
    deceleration \citep[e.g.][]{Matsumoto2013};
  \item \emph{Core-collapse supernovae}, where the shock-heated matter
    decelerates through composition interfaces \citep{Couch2014};
  \item \emph{Gamma-ray bursts} (GRBs), where the relativistic blast
    wave sweeps up ambient medium and the contact discontinuity is
    RT-unstable \citep{Duffell2013}.
  \end{itemize}
\end{enumerate}

In what follows we derive the relativistic perturbation equations for a
stratified fluid and solve the inviscid problem in closed form for two
canonical density profiles.

%% ================================================================
\section{Perturbation equations for a stratified relativistic fluid}
\label{sec:rel-91}

\subsection{Relativistic hydrostatic equilibrium}
\label{sec:rel-91-equil}

Consider a fluid at rest in a uniform gravitational field of strength~$g$
directed along the negative $z$-axis.  In the fluid rest frame the
four-velocity is $u^{\mu} = (c,\,0,\,0,\,0)$ and the energy-momentum
tensor of a perfect fluid takes the form
\begin{equation}
\label{eq:rel-10-1}
T^{\mu\nu} = w\,u^{\mu}\,u^{\nu} + p\,g^{\mu\nu},
\qquad w = \frac{\varepsilon+p}{c^{2}}.
\end{equation}
The relativistic hydrostatic equilibrium condition follows from the
spatial projection of $\nabla_{\mu}T^{\mu\nu}=0$:
\begin{equation}
\label{eq:rel-10-2}
\frac{dp}{dz} = -\,w\,g\,c^{2}\,\frac{1}{c^{2}} = -\,w\,g
\;=\; -\frac{(\varepsilon+p)}{c^{2}}\,g.
\end{equation}
Equation~\eqref{eq:rel-10-2} is the Tolman--Oppenheimer--Volkoff
relation in the weak-field (flat-space) limit.  It reduces to the
Newtonian $dp/dz = -\rho g$ when $\varepsilon \approx \rho c^{2}$ and
$p \ll \varepsilon$.

The enthalpy profile is fixed by a barotropic equation of state
$p = p(\varepsilon)$ together with~\eqref{eq:rel-10-2}; the equilibrium
quantities $\varepsilon_{0}(z)$, $p_{0}(z)$, and $w_{0}(z)$ are thus
functions of the vertical coordinate alone.

\subsection{Eulerian perturbation of the energy-momentum tensor}
\label{sec:rel-91-pert}

We perturb the equilibrium by writing
\begin{equation}
\label{eq:rel-10-3}
\varepsilon = \varepsilon_{0}(z) + \delta\varepsilon, \qquad
p = p_{0}(z) + \delta p, \qquad
u^{i} = \delta u^{i} \quad (i=1,2,3),
\end{equation}
where all perturbation amplitudes are first order in a formal smallness
parameter.  The four-velocity normalization $u^{\mu}u_{\mu}=-c^{2}$
yields, to first order,
\begin{equation}
\label{eq:rel-10-4}
\delta u^{0} = 0.
\end{equation}
Hence the perturbed energy-momentum tensor,
$T^{\mu\nu} = T^{\mu\nu}_{0} + \delta T^{\mu\nu}$, gives
\begin{align}
\label{eq:rel-10-5}
\delta T^{00} &= \delta\varepsilon, \\[4pt]
\label{eq:rel-10-6}
\delta T^{0i} &= (\varepsilon_{0}+p_{0})\,\frac{\delta u^{i}}{c^{2}}
  = w_{0}\,\delta u^{i}, \\[4pt]
\label{eq:rel-10-7}
\delta T^{ij} &= \delta p\;\delta^{ij}.
\end{align}

\subsection{Linearised equations of motion}
\label{sec:rel-91-linear}

Conservation of energy-momentum, $\nabla_{\mu}\delta T^{\mu\nu}=0$,
yields the linearised relativistic Euler equations.  In the weak-field
limit (Minkowski background with an external gravitational potential
$\Phi = gz$) these read:

\medskip\noindent
\textit{Momentum equation} ($\nu = i$):
\begin{equation}
\label{eq:rel-10-8}
w_{0}\,\frac{\partial(\delta u^{i})}{\partial t}
= -\frac{\partial(\delta p)}{\partial x^{i}}
  - g\,\delta w\;\delta^{i3},
\end{equation}
where
\begin{equation}
\label{eq:rel-10-9}
\delta w = \frac{\delta\varepsilon + \delta p}{c^{2}}
\end{equation}
is the Eulerian perturbation of the enthalpy density.

\medskip\noindent
\textit{Energy equation} ($\nu = 0$):
\begin{equation}
\label{eq:rel-10-10}
\frac{\partial(\delta\varepsilon)}{\partial t}
+ (\varepsilon_{0}+p_{0})\,\frac{\partial(\delta u^{i})}{\partial x^{i}}
+ \delta u^{3}\,\frac{d\varepsilon_{0}}{dz}
+ \delta u^{3}\,\frac{dp_{0}}{dz}
= 0.
\end{equation}
Using the equilibrium relation~\eqref{eq:rel-10-2} the last two terms
combine:
\begin{equation}
\label{eq:rel-10-11}
\frac{\partial(\delta\varepsilon)}{\partial t}
+ (\varepsilon_{0}+p_{0})\,\nabla\!\cdot\!\delta\boldsymbol{u}
+ \delta u^{3}\,\frac{d\varepsilon_{0}}{dz}
- w_{0}\,g\,\delta u^{3}
= 0.
\end{equation}

For an \emph{incompressible} relativistic fluid
($\nabla_{\mu}u^{\mu}=0$, i.e.\ $\nabla\!\cdot\!\delta\boldsymbol{u}=0$
to first order), the continuity constraint becomes
\begin{equation}
\label{eq:rel-10-12}
\frac{\partial(\delta u^{1})}{\partial x}
+ \frac{\partial(\delta u^{2})}{\partial y}
+ \frac{\partial(\delta u^{3})}{\partial z} = 0,
\end{equation}
and the energy perturbation is advective:
\begin{equation}
\label{eq:rel-10-13}
\frac{\partial(\delta\varepsilon)}{\partial t}
= -\delta u^{3}\,\frac{d\varepsilon_{0}}{dz}
  + w_{0}\,g\,\delta u^{3}.
\end{equation}

\subsection{Normal-mode reduction}
\label{sec:rel-91-normalmodes}

Seeking solutions proportional to
$\exp(ik_{x}x + ik_{y}y + \sigma t)$, with
$k^{2}=k_{x}^{2}+k_{y}^{2}$ and $D = d/dz$, the momentum
equations~\eqref{eq:rel-10-8} and incompressibility~\eqref{eq:rel-10-12}
combine---exactly as in the Newtonian derivation (cf.\
equations~(10-14)--(10-22))---to give a single equation for the vertical
velocity perturbation $\hat{w}(z)$:
\begin{equation}
\boxed{
D\!\bigl(w_{0}\,D\hat{w}\bigr) - k^{2}\,w_{0}\,\hat{w}
= -\frac{k^{2}}{\sigma^{2}}\,g\,(Dw_{0})\,\hat{w}.
}
\label{eq:rel-10-14}
\end{equation}
This is the \textbf{relativistic Rayleigh--Taylor eigenvalue equation}.
It is identical in form to the classical equation~(10-42) with the
substitution
\begin{equation}
\label{eq:rel-10-15}
\rho \;\longrightarrow\; w_{0}
= \frac{\varepsilon_{0}+p_{0}}{c^{2}}.
\end{equation}
The boundary conditions are
\begin{equation}
\label{eq:rel-10-16}
\hat{w} = 0 \quad\text{on the bounding surfaces.}
\end{equation}

\paragraph{Physical content.}
Equation~\eqref{eq:rel-10-14} embodies the key relativistic modification:
the gravitational weight \emph{and} the inertia of each fluid layer are
governed by $w_{0} = (\varepsilon+p)/c^{2}$, not by $\rho$ alone.  A
layer with large internal (thermal) energy is both heavier and harder to
accelerate; these two effects partially compensate, but generically the
growth rate is altered compared with the Newtonian prediction whenever
$p/\varepsilon$ is non-negligible.

\subsection{Surface tension at interfaces between relativistic fluids}
\label{sec:rel-91-surftens}

When the enthalpy density $w_{0}$ suffers a discontinuity at an
interface $z = z_{s}$, surface-tension forces must be included.
In a relativistic setting the surface energy density (surface tension)
$\mathcal{T}$ contributes to the energy-momentum tensor of the
interface through the Israel junction conditions.  For a weakly curved
interface the net effect on the linearised vertical momentum equation is
the same as in the Newtonian case:
\begin{equation}
\label{eq:rel-10-17}
\Delta_{s}\!\bigl(w_{0}\,D\hat{w}\bigr)
= -\frac{k^{2}}{\sigma^{2}}
  \bigl[g\,\Delta_{s}(w_{0}) - k^{2}\mathcal{T}_{s}\bigr]\,\hat{w}_{s},
\end{equation}
where $\Delta_{s}(f) = f(z_{s}+0)-f(z_{s}-0)$ is the jump at the
interface.  Equation~\eqref{eq:rel-10-17} is the relativistic
generalisation of the classical equation~(10-49).

The relativistic surface energy $\mathcal{T}$ receives contributions from
both the rest-mass surface tension and the thermal/binding energy of the
interface; for a QCD deconfinement front, for example, $\mathcal{T}$ is
of order the bag constant times a characteristic length scale.

%% ================================================================
\section{The inviscid case---relativistic}\label{sec:rel-92}

We specialise equation~\eqref{eq:rel-10-14} to two analytically
tractable configurations.

\subsection{Two uniform relativistic fluids separated by a horizontal
boundary}
\label{sec:rel-92-two}

Let two uniform perfect fluids with constant enthalpy densities
\begin{equation}
\label{eq:rel-10-18}
w_{1} = \frac{\varepsilon_{1}+p_{1}}{c^{2}}
\quad (z<0),
\qquad
w_{2} = \frac{\varepsilon_{2}+p_{2}}{c^{2}}
\quad (z>0)
\end{equation}
be separated by a horizontal interface at $z=0$.  Within each uniform
region $Dw_{0}=0$, so equation~\eqref{eq:rel-10-14} reduces to
\begin{equation}
\label{eq:rel-10-19}
(D^{2}-k^{2})\,\hat{w} = 0,
\end{equation}
with solutions vanishing at $z\to\mp\infty$:
\begin{equation}
\label{eq:rel-10-20}
\hat{w}_{1} = A\,e^{+kz}\;(z<0),
\qquad
\hat{w}_{2} = A\,e^{-kz}\;(z>0).
\end{equation}
Applying the junction condition~\eqref{eq:rel-10-17} at $z=0$:
\begin{equation}
\label{eq:rel-10-21}
-k(w_{2}+w_{1})\,A
= -\frac{k^{2}}{\sigma^{2}}\bigl[g(w_{2}-w_{1})-k^{2}\mathcal{T}\bigr]\,A.
\end{equation}
Hence the \textbf{relativistic Rayleigh--Taylor dispersion relation}:
\begin{equation}
\boxed{
\sigma^{2}
= g\,k\!\left(
  \frac{w_{2}-w_{1}}{w_{2}+w_{1}}
  - \frac{k^{2}\mathcal{T}}{g(w_{2}+w_{1})}
\right)\!.
}
\label{eq:rel-10-22}
\end{equation}
This is the relativistic counterpart of the classical
equation~(10-51).  We define the \textbf{relativistic Atwood number}
\begin{equation}
\label{eq:rel-10-23}
A_{\mathrm{rel}}
= \frac{w_{2}-w_{1}}{w_{2}+w_{1}}
= \frac{(\varepsilon_{2}+p_{2})-(\varepsilon_{1}+p_{1})}
       {(\varepsilon_{2}+p_{2})+(\varepsilon_{1}+p_{1})}.
\end{equation}
In the non-relativistic limit ($p \ll \varepsilon \approx \rho c^{2}$),
$w \to \rho$ and $A_{\mathrm{rel}} \to
(\rho_{2}-\rho_{1})/(\rho_{2}+\rho_{1})$, recovering the classical
Atwood number.

\paragraph{Stability criterion.}
\begin{itemize}
\item If $w_{2} < w_{1}$ (lighter fluid on top), $\sigma^{2}<0$
  for all~$k$ and the interface is \emph{stable} (oscillatory modes).
\item If $w_{2} > w_{1}$ (heavier on top), $\sigma^{2}>0$ for
  $k < k_{c,\mathrm{rel}}$ where
\begin{equation}
\label{eq:rel-10-24}
k_{c,\mathrm{rel}}
= \left[\frac{(w_{2}-w_{1})\,g}{\mathcal{T}}\right]^{1/2}\!.
\end{equation}
Surface tension stabilises wavelengths shorter than
$\lambda_{c}=2\pi/k_{c,\mathrm{rel}}$.
\end{itemize}

\paragraph{Maximum growth rate.}
The wavenumber of fastest growth is
\begin{equation}
\label{eq:rel-10-25}
k_{*} = k_{c,\mathrm{rel}}/\sqrt{3},
\end{equation}
at which
\begin{equation}
\label{eq:rel-10-26}
\sigma_{*}
= \left[
  \frac{2}{3^{3/2}}\,
  \frac{(w_{2}-w_{1})\,g^{2}}{(w_{2}+w_{1})\,\mathcal{T}^{1/2}}
\right]^{1/2}\!,
\end{equation}
reproducing equation~(10-55) with $\rho\to w$.

\paragraph{Role of pressure in inertia.}
A hot, radiation-dominated layer ($p=\varepsilon/3$) has
$w = 4\varepsilon/(3c^{2})$, which exceeds the rest-mass density by a
factor $4/(3\times\rho c^{2}/\varepsilon)$.  Consequently, two fluids
with identical rest-mass densities but different temperatures can exhibit
Rayleigh--Taylor instability in the relativistic regime (whereas they
would be neutrally stable in Newtonian theory).  This is a genuinely
relativistic phenomenon without classical counterpart.

\subsection{Exponentially varying enthalpy density}
\label{sec:rel-92-exp}

Suppose
\begin{equation}
\label{eq:rel-10-27}
w_{0}(z) = w_{00}\,e^{\beta z},
\end{equation}
where $\beta$ is a constant.  Equation~\eqref{eq:rel-10-14} then
becomes
\begin{equation}
\label{eq:rel-10-28}
D^{2}\hat{w} + \beta\,D\hat{w}
- k^{2}\!\left(1 - \frac{g\beta}{\sigma^{2}}\right)\hat{w} = 0.
\end{equation}
This is formally identical to the classical equation~(10-57), with
$\beta$ now being the gradient parameter of $w_{0}$ rather than of
$\rho$.

For a fluid confined between rigid planes at $z=0$ and $z=d$, the
eigenfunction satisfying $\hat{w}(0)=\hat{w}(d)=0$ is
\begin{equation}
\label{eq:rel-10-29}
\hat{w}(z) = C\,e^{-\beta z/2}\,\sin\!\left(\frac{m\pi z}{d}\right),
\qquad m = 1, 2, 3, \ldots
\end{equation}
and the dispersion relation is
\begin{equation}
\label{eq:rel-10-30}
\frac{g\beta}{\sigma^{2}}
= 1 + \frac{\tfrac{1}{4}\beta^{2}d^{2}+m^{2}\pi^{2}}{k^{2}d^{2}}.
\end{equation}

\paragraph{Stability.}
If $\beta<0$ (enthalpy decreasing upward), $\sigma^{2}<0$ and the
stratification is stable; if $\beta>0$ (enthalpy increasing upward),
$\sigma^{2}>0$ and the stratification is unstable.  The relativistic
Atwood number for the exponential profile is
\begin{equation}
\label{eq:rel-10-31}
A_{\mathrm{rel}}(z)
= \frac{w_{0}(z+\delta z)-w_{0}(z-\delta z)}
       {w_{0}(z+\delta z)+w_{0}(z-\delta z)}
\;\xrightarrow{\delta z\to 0}\;
\tanh\!\left(\tfrac{1}{2}\beta\,\delta z\right)
\approx \tfrac{1}{2}\beta\,\delta z,
\end{equation}
confirming the identification $\beta \leftrightarrow D(\ln w_{0})$ as the
local measure of stratification.

\subsection{Causality check}
\label{sec:rel-92-causality}

The Rayleigh--Taylor instability is a purely growing (or decaying) mode:
$\sigma$ is real.  Nonetheless, the perturbation equations must be
checked for causal consistency.

\begin{enumerate}
\item \textbf{Phase velocity of RT modes.}  The normal-mode ansatz
  $\sim\exp(ik_{x}x+ik_{y}y+\sigma t)$ with real $\sigma$ describes a
  mode that is spatially periodic but temporally exponential.  There is
  no phase velocity to compare with~$c$: the ``perturbation'' does not
  propagate horizontally---it grows (or decays) in place.  Causality is
  therefore not violated by the RT mode itself.

\item \textbf{Sound-speed bound.}  The derivation assumed an
  incompressible fluid.  In a compressible relativistic fluid, RT modes
  with wavelengths shorter than the Jeans-like scale
  $\lambda_{\mathrm{J}} \sim c_{\mathrm{s}}/\sqrt{g\,|Dw_{0}|/w_{0}}$
  couple to sound waves; the sound speed must satisfy
  $c_{\mathrm{s}}^{2} = (\partial p/\partial\varepsilon)_{s} \le c^{2}$
  for causal thermodynamics.

\item \textbf{Dissipative corrections.}  Viscosity and heat conduction,
  when included via the BDNK first-order causal formalism
  \citep{Bemfica2018,Kovtun2019}, produce constitutive relations
  $\Pi^{\mu\nu}=-2\eta\,\sigma^{\mu\nu}$ with transport coefficients
  constrained by frame-dependent causality bounds, guaranteeing that
  even the short-wavelength cutoff from dissipation respects
  causality---unlike a naive relativistic Navier--Stokes extension.
  No relaxation times $\tau_\pi$ or $\tau_q$ are needed; causality is
  ensured by the BDNK frame coefficients.
\end{enumerate}

We conclude that the relativistic RT growth rate~\eqref{eq:rel-10-22} is
physically well-posed and causal for all wavenumbers within the regime
of validity of the incompressible, inviscid approximation.

%% ================================================================

\begin{center}
\rule{0.5\textwidth}{0.4pt}
\end{center}

\noindent\textbf{Summary of relativistic modifications
(Chapter~X, \S\S\,90--92).}
\begin{enumerate}
\item The eigenvalue equation~\eqref{eq:rel-10-14} has the same
  Sturm--Liouville structure as its Newtonian counterpart, with $\rho$
  replaced everywhere by $w=(\varepsilon+p)/c^{2}$.
\item The relativistic Atwood number is
  $A_{\mathrm{rel}}=(w_{2}-w_{1})/(w_{2}+w_{1})$.
\item Pressure contributes to both weight and inertia; two fluids of
  equal rest-mass density but different internal energies can be
  RT-unstable.
\item The growth rate $\sigma$ is real; causality is satisfied
  automatically for the incompressible limit and is enforced by the
  BDNK first-order causal formalism for dissipative corrections.
\end{enumerate}

%% ================================================================
\section{Astrophysical application: Shock-driven Rayleigh--Taylor
instability in core-collapse supernovae}\label{sec:rel-92-astro}

A prime astrophysical setting for the relativistic Rayleigh--Taylor
instability is the \emph{core-collapse supernova} (CCSN).  During
the post-bounce phase, the stalled accretion shock decelerates
as it encounters composition interfaces (Si/O, O/He), creating
conditions for RT growth at those boundaries
\citep{Couch2015,Mueller2012}.  The post-shock matter is heated to
temperatures $T \sim 5$--$10\;\mathrm{MeV}$, so that the thermal
pressure contribution to the enthalpy density is non-negligible:
$p/(\varepsilon) \sim 0.05$--$0.20$.  Consequently, the
\emph{relativistic Atwood number}
\begin{equation}
A_{\mathrm{rel}}
= \frac{w_{2}-w_{1}}{w_{2}+w_{1}}
= \frac{(\varepsilon_{2}+p_{2})-(\varepsilon_{1}+p_{1})}
       {(\varepsilon_{2}+p_{2})+(\varepsilon_{1}+p_{1})}
\end{equation}
differs from the classical value
$A = (\rho_{2}-\rho_{1})/(\rho_{2}+\rho_{1})$.  For hot post-shock
matter with $\xi = p/(\rho c^{2}) \sim 0.1$--$0.2$, the relativistic
correction modifies the growth rate by $\sim 5$--$15\%$, which is
significant for predicting the degree of mixing and the ultimate
success or failure of the explosion mechanism
\citep{Burrows2013,Ott2013}.

Figure~\ref{fig:supernova-RT} illustrates this effect by plotting the
RT growth rate as a function of the density ratio across the
composition interface for several values of the relativistic
parameter~$\xi$.  The right panel compares the classical and
relativistic Atwood numbers.

\begin{figure}[htbp]
\centering
\includegraphics[width=0.95\textwidth]{fig_supernova_RT.pdf}
\caption{Shock-driven Rayleigh--Taylor instability in core-collapse
  supernovae.  \textit{Left:} Growth rate $\sigma$ as a function of
  the density ratio $\rho_{2}/\rho_{1}$ across the composition
  interface, for classical ($\xi=0$) and several relativistic values
  of $\xi = p/(\rho c^{2})$.  \textit{Right:} Classical vs.\
  relativistic Atwood number.  Parameters correspond to a proto-neutron
  star of mass $1.4\,M_{\odot}$ and shock radius $150\;\mathrm{km}$.}
\label{fig:supernova-RT}
\end{figure}
